\documentclass[activite]{thedocument}
\usepackage{tkz-euclide}
\usepackage{enumitem}
\startdatakeys
\source{}
\cycle{}
\competencesDuSocle{}
\connaissancesRequises{}
\competencesTravaillees{}
\enddatakeys
\begin{document}
    {\centering \begin{tikzpicture}\sf
        \tkzDefPoint(0,0){O}
        \tkzDefPoint(1,0){A}
        \tkzDefPoint(-1,0){-A}
        \tkzDefPoint(2,0){B}
        \tkzDefPoint(-2,0){-B}
        \tkzDefPoint(3,0){C}
        \tkzDefPoint(-3,0){-C}
        \tkzDefPoint(4,0){D}
        \tkzDefPoint(-4,0){-D}
        \tkzDefPoint(2.5,0){E}
        \tkzDefPoint(5,0){O1}
        \tkzDefPoint(-5,0){O2}
        \draw[-latex](O2)--(O1);
        \draw[dotted](-6,-3)rectangle(6,3);
        \foreach \x in {-5,...,4} {
            \draw[line width=1pt] (\x,0.1cm) -- (\x,-0.1cm);
            \draw (\x+0.5,0.05cm)--(\x+0.5,-0.05cm);
            }
        \draw node[below] at(O){0\strut};
        \draw node[below] at(A){1\strut};
        \draw node[below] at(-A){-1\strut};
    \end{tikzpicture}}\vskip20pt
    Pour calculer avec les nombres relatifs, on considère que~:
    \begin{itemize}[label=$-$]
        \item les nombres {\bfseries positifs} représentent \doted{100pt}.
        \item les nombres {\bfseries négatifs} représentent \doted{100pt}.
    \end{itemize}
    \noindent{\bfseries Exemple~:}
    \begin{itemize}[label=$-$]
        \item Un gain de 37 euros est représenté par le nombre \doted{40pt}.
        \item Un perte de 12 euros est représenté par le nombre \doted{40pt}.
    \end{itemize}
    \noindent{\bfseries Exemple~:}\vskip10pt
    \hbox{\hbox to.5\linewidth{$(+12)+(+9)=$\hfil}\hbox to.5\linewidth{$(-7)+(-16)=$\hfil}}\vskip15pt
    \hbox{\hbox to.5\linewidth{$(+30)+(-5)=$\hfil}\hbox to.5\linewidth{$(+12)+(-18)=$\hfil}}\vskip15pt
    \hbox{\hbox to.5\linewidth{$(-4)+(+7)=$\hfil}\hbox to.5\linewidth{$(-10)+(+9)=$\hfil}}
\end{document}